%
%
The past decade of deep learning has produced systems of impressive capability, 
yet their deployment in demanding real-world settings consistently exposes a 
recurring cluster of failures: brittleness to distributional shift, inability 
to compose learned concepts in novel ways, and error accumulation that renders 
long-horizon planning unreliable. This thesis argues that these failures are 
not incidental engineering shortcomings but structural consequences of how 
current architectures represent knowledge. We propose a research direction that 
moves from diagnosing these failures, through mitigating them with principled 
constraints, toward a paradigm shift grounded in neuro-inspired computation.

We begin with a comprehensive survey and taxonomy of agents for computer use 
(ACUs), systems that automate complex tasks on digital devices given 
natural-language instruction. Covering 87 research papers and 33 benchmark 
datasets, we identify six persistent research gaps, including insufficient 
generalisation, limited planning, and a disconnect between laboratory evaluation 
and real-world deployment. We argue that these gaps are symptomatic of 
structural limitations in the underlying learning paradigm rather than a matter 
of scale alone. Two supporting studies corroborate this diagnosis from 
independent directions. A benchmark of efficient large language models shows 
that the compression methods currently needed to make these models accessible 
systematically degrade the reasoning capabilities they are meant to preserve. 
A compositional visual reasoning framework further demonstrates that 
state-of-the-art architectures, including those with strong inductive biases 
for compositionality, consistently fail to generalise to novel rule combinations.

To understand the origin of these failures, we develop a mathematical framework 
for analysing the representational geometry of self-attention matrices. We prove 
that the choice of training objective, bidirectional versus autoregressive, 
dictates a precise structural signature in the learned weight matrices: symmetry 
in the former, directionality and column dominance in the latter. These results 
are validated across ModernBERT, GPT, LLaMA3, and Mistral on text, vision, and 
audio modalities. They establish that the geometry of learned representations is 
a mathematical consequence of the training objective. Current architectures do 
not discover structure in the world; they embed the statistical structure of 
their loss function.

Given this diagnosis, we investigate whether imposing hard, mathematically 
grounded constraints on deep learning systems can correct specific failure modes. 
We introduce \textit{Koopman-JEPA}, a self-supervised world model that lifts 
complex visual observations into a latent space governed by globally linear 
dynamics. In standard nonlinear world models, approximation errors compound 
exponentially across rollout steps via the Lipschitz constant of the composed 
functions. By replacing these nonlinear rollouts with matrix multiplications 
over a stable linear operator, error accumulation is bounded to grow linearly 
with horizon. Three variants with progressively stronger spectral constraints 
on the transition matrix culminate in a Cayley parameterisation that 
architecturally enforces strict stability and satisfies the stabilisability 
precondition of discrete-time linear quadratic regulation. Koopman variants 
reduce long-horizon prediction error by up to 99\% on action-free tasks and 
91\% on action-conditioned tasks at a ten-step horizon, while improving sample 
efficiency by approximately sixteen-fold when only ten percent of training data 
is available. Two complementary studies show that the constraint principle 
generalises across domains. Encoding tissue-specific Hounsfield unit ranges as 
binary input masks compensates for data scarcity in clinical CT segmentation, 
and a hybrid retrieval architecture that grounds predictions via external 
structured knowledge achieves first place on the CLEF CheckThat! 2025 
development leaderboard without any external training data.

Finally, we ask whether constraint engineering alone is sufficient or whether 
the representational substrate must itself change. We introduce the 
\textit{Cooperative Network Architecture} (CNA), which replaces distributed 
activation patterns with dynamically assembled, recurrently connected networks 
of neurons called \textit{nets}, built from overlapping fragments learned 
without supervision from regularities in sensory input. Because nets are 
composable by construction, CNA exhibits figure completion, noise resilience, 
and generalisation to out-of-distribution patterns without retraining. These 
properties emerge from the representational principle of the architecture rather 
than from loss engineering.

Together, these contributions trace a path from empirical diagnosis through 
principled mitigation to architectural reconception. We conclude that 
constraint-based approaches offer the most immediate practical gains, but that 
achieving robust and composable intelligence ultimately requires representational 
structures that embed invariance and compositionality by design, a direction in 
which neuro-inspired computation remains the most principled guide.
